% The opening sections, shared by the rulebook and the platform manuals (docs/manual/*).
% Each document sets its own heading above this.
You are a \term{\V{playertitle}}. You \term{\V{jackin}} to the region's \term{\V{grid}}, examine your \term{\V{kit}} and search for unpatched \term{\Vp{node}} to \term{\V{spike}}.  When two \Vp{spike} are linked, you've gained access. Syndicates want control of the \V{grid} and will pay a \term{\V{finderscut}} in their \term{\V{coin}} for this access. This becomes their \term{\V{hive}}. \VCp{hive} spread across the \V{grid} with each new \V{spike}. As a \V{hive} grows, so does its \V{coin}'s value. 

But syndicates do not share. When two \Vp{hive} collide, the bigger one \term{\Vp{hivejack}} the smaller, absorbing it. The smaller \V{hive}'s top two \V{coin} backers make bank. Hope that you're one of them.

Eventually, the \Vp{hive} grow too big to remain unnoticed. The carriers catch on and \term{\V{nullroute}} the entire region. Your gig is up. When the routes go dark, the \V{player} holding the most \Game{currency} wins.

You've learned the hard way: the real \Game{currency} is not in brokering \V{spike}s, but in playing the syndicates against each other. 

\subsection{Game Description}
SYN-ACK-WIRE is a strategic, turn-based, online multiplayer game for FujiNet-enabled computers set in a cyberpunk world. The table seats two to six players, including bots, so you can play solo.  Take turns, spiking the network and dabbling in \V{coin}, until an end-of-game condition exists.

\subsection{Object of the game}
Have the most \V{money} (\Game{currency}) when the region is \V{nullroute}d.

\subsection{End Game Conditions}
Play can end when:
\begin{enumerate}
  \item Any hive has become a \V{totalcompromise} by spreading to 41+ nodes.
  \item All hives have become immune to \Vp{hivejack} (\term{\V{safe}}) by spreading to 11+ nodes. (automatic)
\end{enumerate}
